\documentclass[10pt,a4paper]{article}
\usepackage[T1]{fontenc}
\usepackage{lmodern,amsmath,amssymb,microtype,booktabs,graphicx}
\usepackage[margin=20mm]{geometry}
\usepackage[hidelinks]{hyperref}
\pdfinfoomitdate=1
\pdftrailerid{}
\pdfsuppressptexinfo=15
\setlength{\parindent}{0pt}
\setlength{\parskip}{5pt}
\begin{document}
\begin{center}
{\Large Candidate causes of the phase-local H$\alpha$ contrast}\par\vspace{4pt}
\texttt{exojax-escape} --- 27 November 2026
\end{center}

\textbf{Question and limits.}
The first night's twelve exposures at $\phi=0.3761496$--$0.3932226$ have the
most consistent complete-band residual signs: the blue peak is too weak,
and the modeled core and red band are too bright. The primary blue/core/red
residuals are $(+591,-585,-216)$ ppm; complementary pixels give
$(+402,-449,-223)$ ppm. These are processed native log-contrast means, not
absolute planetary fluxes. The blue-minus-core residual is $1176/852$ ppm;
its conditional errors are $355/345$ ppm, while empirical time-block
proxies span $544$--$1052$ and $386$--$901$ ppm. The short interval and unknown
covariance still prevent a calibrated physical detection.

Three parallel numerical investigations retain the same native sampling,
stellar mask, full-data reference, fixed noise hyperparameters and
complete-band maps. Conditional GP means are recomputed when a prediction
changes. Complete blue support ends at $\phi=0.39322$ and is absent on the
second night. No missing band is extrapolated. The selected interval is a
post-hoc diagnostic target, not an independent validation sample.

\textbf{Executed comparison.} Positive $\Delta Q=Q_{\rm base}-Q_{\rm case}$ denotes a better full-native-data score. Each row below refits the same two global RV/amplitude parameters; the alternative pixels also use their own refitted baseline. Noise hyperparameters are fixed. Values are rounded ppm, and the original fixed-baseline alternative residuals in the introduction therefore differ slightly from this table.

\begin{center}\small\begin{tabular}{lrrr}\toprule

Shape & Primary B/C/R & Alternate B/C/R & $\Delta Q$, primary/alternate\\\midrule

Baseline & +591/-585/-216 & +424/-463/-183 & -0.00/-0.00\\

Upper source $\times$0.5 & +444/-221/-294 & +280/-121/-245 & +23.73/+19.59\\

Inner $\times$1.5, upper $\times$0.5 & +416/-262/-264 & +260/-168/-213 & +27.28/+19.01\\

Upper opacity $\times$2.014, fixed emission & +464/-168/-330 & +277/-63/-273 & +13.15/+11.81\\

Free model 1, original shape & +184/-165/-184 & +36/-82/-117 & +21.83/-0.52\\

Free model 1, refitted shape & +168/-371/-113 & +33/-282/-42 & +37.75/+2.46\\

Fitted smooth longitude map & +578/-592/-200 & +414/-471/-168 & +0.46/-0.97\\

Localized source patch $\times$4 & +476/-606/-15 & +329/-494/+12 & -13.58/-22.49\\

Upper LOS offset +5 km/s & +362/-565/-99 & +240/-473/-49 & +20.40/-7.83\\

Upper radial speed +13 km/s & +864/-623/-162 & +680/-512/-157 & -72.83/-59.08\\

\bottomrule\end{tabular}\end{center}

\textbf{Source and opacity.} Reducing upper emission improves the trough, but appreciable blue residual remains. Doubling upper opacity at fixed source gives only a small response; increasing the lower-state reservoir while keeping emission approximately fixed behaves much like lowering the source. A stronger inner source also adds red emission. Matching only the selected blue/core means with two source coefficients requires large changes and creates substantial collateral errors; it is not an admissible all-data solution.

\textbf{Kinematics and longitude.} A smooth positive longitude map improves primary $Q$ by less than one. A localized bright region can increase blue emission but barely deepens the core and worsens the global comparison. A uniform upper LOS offset is not a radial wind: the actual ray-projected radial acceleration moves the blue/core residual in the wrong direction. The LOS offset's global preference also changes with pixel selection.

\textbf{Stellar/calibration controls.} With the original planetary RV/amplitude held fixed, the full-data stellar-width drift changes the first night's width by at most 0.64\%; its exact target residuals are approximately $(+551,-548,-185)$ ppm and its primary $\Delta Q$ is 22.18. This removes only about 6.5\% of the blue-minus-core discrepancy. Depth pulsation and broad-curvature drift remove less than 1\%. A target-selected 8.45\% width drift can mimic much of the local shape, but its exact full-data $Q$ worsens by about 1471. It is not a viable correction supported by these data.

\textbf{Model 1 is informative, but not a physical identification.} Its original shape reduces both target residuals much more than uniform source scaling. Reoptimizing its shape further improves global $Q$ on primary pixels while leaving a larger local core error. The alternative-pixel global gain is small or absent. Thus a good local shape, a better global score and a physically realizable atmosphere are distinct requirements.


\clearpage
\textbf{What a physical explanation would have to do.}
For an illustrative emitting inner layer seen through an upper layer,
\[
 I_\nu=I_{\nu,\mathrm{in}}e^{-\tau_\nu}
       +S_{\nu,\mathrm{up}}(1-e^{-\tau_\nu}).
\]
Lowering the upper source function can deepen an optically thick core.
Increasing optical depth alone saturates once $\tau_\nu\gtrsim1$ and does
not remove the upper layer's own emission. At fixed lower-level population,
reducing the upper-level population lowers $S_\nu$ with only a small
stimulated-emission correction to opacity. H$\alpha$ therefore constrains
both the $n=3$ to $n=2$ population ratio and the distribution of $n=2$ in
position and velocity. It does not supply a free, physically independent
source-function parameter.

A partially covering, dim layer changes an intrinsic spectrum by
$\Delta I_\nu=f(S_\nu-I_{\nu,\mathrm{in}})(1-e^{-\tau_\nu})$.
It can suppress the core where $S_\nu<I_{\nu,\mathrm{in}}$. In the pure-absorption
limit ($S_\nu\simeq0$), it cannot create intrinsic blue emission.
The observed target additionally passes through
common-reference subtraction and GP conditioning; its signed band responses
must be evaluated with that operator, as above.

\textbf{The free model points to the radial placement of emission and absorption.}
Model 1 (original shape) is not just the physical atmosphere with a lower upper source
function. Above $1.3R_p$ it has an $n=2$ column of
$4.08\times10^{13}$ rather than $5.85\times10^{12}\,\mathrm{cm}^{-2}$;
the static nominal-center optical depths are about 15.8 and 3.10.
Above $1.725R_p$, the contrast is much larger: model 1 has
$\tau\simeq4.88$, while the physical atmosphere has $\tau\simeq0.0456$.
In this static radial diagnostic the outward $\tau=1$ radius is
$2.25R_p$ versus $1.59R_p$; these are not inferred observer formation radii.
At $1.94R_p$ it retains $n_2\simeq1401$ rather than
$13\,\mathrm{cm}^{-3}$. Its source is bright around $1.5R_p$, then falls
by more than an order of magnitude while appreciable absorbing material
remains farther out. Its gas is hotter, not cooler, than the physical
upper atmosphere: it does not demonstrate a solution from thermal narrowing.

The free model's total $n=2$ column is actually smaller. Thus the useful
clue is where the absorbing column and source-function transition lie,
not simply adding more absorption everywhere. Plane versus spherical
geometry, off-limb emission, continuum, atomic branches, velocity and
independently prescribed populations also differ, so this comparison
does not isolate a single cause. A matched-geometry transfer comparison
is needed before attributing its advantage to the extended upper column.

\textbf{Priority physical hypotheses: an extended dim absorbing region and nonuniform excitation.}
A dim upper layer in front of bright, approaching material is a possible
realization of the required contrast. Day/night or longitude-dependent
irradiation, cooling, density and line-of-sight coverage can give different
source functions to the emitter and absorber. A layer that sees a smaller
solid angle of bright inner emission could also receive less line pumping.
The tested longitude source maps change emissivity at fixed opacity;
they do not test this coupled spatial variation. The calculation must
preserve excitation, photon and energy budgets: assigning a dark screen
without allowing absorbed photons to repopulate $n=3$ is not a solution.

This interpretation is motivated, not established, by three-dimensional
excited-state escape models
\href{https://doi.org/10.1051/0004-6361/202453309}{(Shaikhislamov et al. 2025)}.
Transit-based asymmetric-tail fits require sustained excited hydrogen far
from the planet and admit stellar alternatives
\href{https://arxiv.org/abs/2301.10406}{(Lowson et al. 2023)}; they are not
evidence for a tail at the present emission phases.

\textbf{A representation limitation must be separated from new physics.}
The physical atmosphere has eight radial cells. The main inner emitting
cell spans $1.2167$--$1.3R_p$ and its upper neighbor $1.3$--$1.725R_p$.
Their rest-frame radial peak line optical depths are about $54$ and $3$;
the next cell is only about $0.04$. The core thus probes a steep transition
represented by one broad upper cell. More observer rays cannot establish
convergence of its temperature and population stratification. No matched,
fully re-solved radial refinement of this final state has established that
adequacy. Atomic truncation ($n_{\max}=15$, resolved through $n=3$, no top-off)
also remains a qualification limit. These are reasons to test the closure,
not evidence that numerical refinement will improve the fit.

\clearpage
\textbf{Additional hypotheses and why they rank below a demonstrated fix.}
\emph{Directional irradiation and time dependence.}
The adopted PHOENIX boundary is fixed and spatially uniform. A gravity-darkened
star on this misaligned orbit gives wavelength-dependent illumination twice
per orbit; published estimates are roughly 10\% in the UV and 1--2\% in the
optical \href{https://arxiv.org/html/2004.14812v1}{(Ahlers et al. 2020)}.
As an explicitly illustrative calculation, a 10\% peak-to-peak sinusoid
changes by at most 1.07\% across this 36.39-minute interval. A large change
of source function would require a nonlinear or spatial response; the
irradiation modulation alone is not a fitted explanation.

The viewing direction changes by $6.15^\circ$ across those exposures, so a
persistent structure can change visibility without being created in 36
minutes. Saved base hydrogen coefficients give gross recombination scales
of about 25 s and 283 s in the two principal cells, versus about 40 minutes
to 3.6 hours farther out. Those historical base rates omit later receiving
PI/RR packets and do not include photon recycling as a relaxation problem.
The upper cell's radial crossing time is about 21 hours. These scales
motivate a coupled time-dependent test; none demonstrates non-equilibrium
ionization or excludes it.

\emph{Frequency redistribution and atomic kinetics.}
Lyman-line transfer can change the $n=2$ reservoir, and frequency-dependent
pumping can change $n=3/n=2$. This warrants an audit of the adopted
redistribution/escape approximation. It does not justify assuming that
direct partial redistribution in H$\alpha$ fixes the trough: solar
calculations can synthesize H$\alpha$ in complete redistribution under
specified Lyman-profile and atmospheric-history conditions
\href{https://arxiv.org/abs/1202.1926}{(Leenaarts et al. 2012)}. Those solar
conditions are not a validation of the planetary closure.

\emph{Already included channels.}
Excited-state photoionization, recombination and Balmer-continuum heating
are present; the physical relevance of the latter is established in
\href{https://arxiv.org/abs/1910.00267}{Garc\'ia Mu\~noz \& Schneider (2019)}.
The previous full spherical H$\alpha$ coupling increased the upper source
function and worsened the fit. In the corrected conditional
H$\alpha$/Ly$\beta$/O I network, deleting the Ly$\beta$ branch changes the
outer H$\alpha$ survival probability by only 0.0534 percentage point.
Other-shell and ionization exits are terminal only for that tracked network,
not permanent photon destruction in the complete atmosphere. These effects should not be
added again or described as untested missing absorbers.

\emph{Stellar and reduction effects.}
Known stellar pulsations motivate time-dependent controls
\href{https://doi.org/10.1051/0004-6361/202243823}{(Jones et al. 2022)}.
The tested line-depth/width/shift and smooth-curvature families quantify
specific alternatives, not every nonradial pulsation pattern, telluric
residual, line-spread-function error or extraction covariance. Their
failure to remove the contrast does not establish an atmospheric origin.

\textbf{Next discriminating calculation.}
First place the free and physical radial source/opacity profiles in the
same observer geometry, swapping only the upper column or the location of
the source transition. This separates the strongest structural clue from
the free model's other freedoms. A physical follow-up must explain how
an extended $n=2$ reservoir survives while $n=3/n=2$ stays small; density
and flow structure, irradiation and shielding are candidate controls.
A minimal two-sector equilibrium is one way to test nonuniform excitation
and coverage. Compare all native phases and both pixel selections; reject a candidate that improves
the selected B/C pair by creating excessive red emission or worsening the
remaining data. Qualify the steep radial transition separately. New
repeated-phase observations or pre-master spectra are still needed to
separate this physical hypothesis from stellar/reduction variability.

\textbf{Reproducibility and MCEG scope.}
The three diagnostic scripts, source/input hashes, numerical controls,
primary-literature notes and compact results are in
\texttt{mceg/evidence/halpha\_causes\_20261127/}.
Large native arrays remain in the recorded local cache with hashes.
No production atmosphere or likelihood is changed. MCEG distinguishes
the finite numerical responses from the unknown physical attribution;
conditional score improvements are neither Bayes factors nor calibrated
detections. Previous reports and evidence remain unchanged.
\end{document}
